% ==========================================================================
% error_book_template —— 错题整理版式(每题含"我的备注"块)
% 由 pdf_gen.py 以双圆括号占位符注入内容后编译(xelatex,需 Noto CJK 字体)。
% 备注块由 pdf_gen.py 在注入题目区时按题渲染。
% 占位符键:TITLE / SUBTITLE / EXAM_DATE / SUBJECT / DURATION /
%           TOTAL_SCORE / NOTICE / QUESTIONS
% ==========================================================================
\documentclass[12pt, a4paper, fontset=none]{ctexart}

\usepackage[top=2.5cm, bottom=2.5cm, left=2.3cm, right=2.3cm]{geometry}
\usepackage{amsmath}
\usepackage{amssymb}
\usepackage{xcolor}
\usepackage{fancyhdr}

% 中文字体(假定系统已安装 Noto CJK 系列)
\setCJKmainfont{Noto Serif CJK SC}
\setCJKsansfont{Noto Sans CJK SC}
\setCJKmonofont{Noto Sans Mono CJK SC}

% 题号标题样式
\ctexset{
  section = {
    format     = {\large\bfseries},
    beforeskip = {1.3em},
    afterskip  = {0.6em},
  },
}

% 页眉页脚
\setlength{\headheight}{16pt}
\pagestyle{fancy}
\fancyhf{}
\fancyhead[L]{\small 错题整理}
\fancyhead[R]{\small ((TITLE))}
\fancyfoot[C]{\small 第~\thepage~页}
\renewcommand{\headrulewidth}{0.4pt}

\begin{document}

% ------------------------------ 标题区 ------------------------------
\begin{center}
  {\huge\bfseries ((TITLE))\par}
  \vspace{0.5em}
  {\large ((SUBTITLE))\par}
  \vspace{0.8em}
  {\normalsize 整理日期:((EXAM_DATE)) \qquad 科目:((SUBJECT)) \qquad
   建议用时:((DURATION))~分钟 \qquad 总分:((TOTAL_SCORE))~分\par}
\end{center}

\vspace{0.4em}

% ------------------------------ 使用说明框 ------------------------------
\noindent\fbox{\parbox{\dimexpr\linewidth-2\fboxsep-2\fboxrule\relax}{%
  \textbf{使用说明}\par\smallskip
  ((NOTICE))\par
}}

\vspace{1em}

% ------------------------------ 错题区(每题含"我的备注"块) ------------------------------
((QUESTIONS))

\vspace{2em}
\begin{center}
  {\small --- 错题整理完 ---}
\end{center}

\end{document}
